\section{Recommendations}\label{recommendations}

\subsection{Recommendations for Valuation Practice}\label{recommendations-for-valuation-practice}

The price surface works best as a triangulation tool rather than a replacement for existing references. The deployed Random Forest was the strongest predictive option in the study, but Chapter 7 showed that the held-out error stayed substantial in some locations and property types, well above the accuracy expected of assessment-grade valuation systems. Practitioners should therefore read the surface alongside BIR zonal values and property-specific comparable evidence instead of treating the model output as a stand-alone final appraisal. When weighing those comparables, the variables that carried the most predictive weight deserve the most attention: Chapters 7 and 8 showed that polycentric distance and property type mattered more consistently than the accessibility composite, so corridor position, proximity to active subcenters, and housing submarket class should come before broad amenity counts.

The spatial sorting result from the Stage 1 OLS model needs careful interpretation. The negative coefficients on security, tourism, and retail density did not mean these uses are always harmful to residential value. They pointed to context effects such as congestion, commercial intensity, or neighborhood sorting. Practitioners should avoid assuming that physical proximity to police infrastructure, tourism strips, or dense commercial activity automatically adds residential value without checking the surrounding land-use context.

The valuation gap between market listings and BIR zonal values should be read as a research signal, not as an appraisal input. The thesis quantified the gap and preserved it as a diagnostic field: for the clean land-to-land comparison on vacant lots, market prices ran roughly three times the BIR zonal benchmark, and listings exceeded the benchmark in close to all LGUs. This indicates that official benchmarks lag the open market across Metro Cebu. It does not, however, establish a correction factor that can be applied during an appraisal. The gap is computed from asking prices through a model that is not assessment-grade, so it points to where benchmarks may be stale and warrants closer study, rather than supplying a value adjustment that practitioners can use directly. Validating the gap against recorded transactions is a prerequisite before any operational use.

\subsection{Recommendations for Policy}\label{recommendations-for-policy}

The findings support continued polycentric planning in Metro Cebu. Chapters 7 and 8 showed that the market did not behave as if Cebu Business Park alone structured residential value; distance signals from Mactan, Consolacion, and Naga also mattered strongly. For regional planners and local governments, investment outside the historical core should keep recognizing the functional role of secondary and corridor-based subcenters rather than concentrating policy attention on the CBP corridor alone.

MCRAI may be useful as a reusable template for parcel-level accessibility, but as a structure to adapt rather than a finished index to copy. The thesis built MCRAI for residential valuation, not for nationwide equity mapping and not for liquidation pricing. Its category-specific radii, locally defined accessibility categories, and positive-only composite rule give a workable starting point for LGU-level planning or assessment work. Its parameters, however, were set on judgment-based grounds rather than estimated from local data, so they should be refined along the lines set out below before the framework is relied on for planning decisions.

\subsection{Recommendations for Future Research}\label{recommendations-for-future-research}

The most important next step is to put the geospatial accessibility features on a more objective and empirical footing. The current Metro Cebu Residential Accessibility Index rests on several design choices that were defensible but largely a matter of judgment rather than estimates from local evidence: the gravity decay parameter was fixed at $\beta = 2.0$, the search radii were set by category, and the composite weights were read from OLS coefficient signs. The framework also adapts an established accessibility-index logic rather than one calibrated specifically to Metro Cebu. Because the MCRAI categories contributed only modestly for condominiums and houses, and because the model's accuracy was limited overall, these arbitrary choices are a real candidate explanation that future work should test directly. A more objective approach would estimate the decay rate and the category weights from the data, compare alternative accessibility formulations against the current one, and treat known negative-effect amenities such as gasoline stations and dense tourism strips as nonlinear or threshold features rather than folding them into a single positive-only score. The aim is a spatial feature set whose construction is reproducible and data-driven rather than dependent on parameter values chosen in advance.

Spatial heterogeneity also deserves more direct estimation. The present models captured non-linearity but did not estimate coefficients that vary across space. Geographically Weighted Regression or Mixed Geographically Weighted Regression is the logical extension, and a follow-up study should test whether the capitalization of Cebu Business Park, Mactan, Consolacion, SRP, and Naga changes materially across submarkets instead of assuming one average distance effect for the whole region.

Broadening the model's inputs is the other major direction, since this study found that geospatial features and open-market listings are not sufficient on their own to value every property type well. Future work should collect richer parcel and structural attributes, such as road frontage and a corner-lot flag, which the vacant-lot data ceiling indicated are priced by the market but are rarely recorded in listings. It should also strengthen the time dimension of the analytical base table by preserving collection dates and adding macroeconomic indicators, such as exchange rates, inflation, and interest rates, so that market cycles and price drift can be modeled directly rather than left out of a cross-sectional design \parencite{udomsap2020}. A further extension is to triangulate predictions against complementary reference prices, such as BIR zonal values and bank reserve prices, instead of relying on a single market-facing target.

\subsection{A Note on Deployment}\label{recommendations-for-deployment}

Deployment was not a primary aim of this study, so only a few practical notes apply to the prototype web application. It should be run inside the project virtual environment so that its dependencies, including SHAP, load correctly, and it should display its open-market-only scope clearly so the price surface is not read as a universal market value. If the prototype is maintained, the point-of-interest inventory, listing coverage, and benchmark schedules will drift over time, so the data and predicted layers should be refreshed and re-verified periodically.

\subsection{Closing Note}\label{closing-note}

This thesis should be read as a demonstrated workflow more than as a finished valuation product. Its main contribution is that the process is reproducible: a custom MCRAI design, a polycentric distance feature set, and a Random Forest plus SHAP pipeline can be rebuilt, checked, and adapted for other Philippine secondary cities with similar data conditions. What this study finished was a defensible template for local valuation analytics, not the last version that Metro Cebu will ever need.
